\subsection{单模}

回忆下面的定义（II, §1, No.~10, 第 212 页）。

定义 1. —— 设 A 是环。若一个 A-模 M 非零且没有异于 0 与 M 的子模，则称它为单模。

A-模 M 是单模当且仅当 M 是其位似环 \(A_{M}\) 上的单模。每个单模都不可分解、长度为 1，从而是原始模（VIII，第 31 页，命题 4, b)）。

例. —— 1) 模 \(A_{s}\) 是单 A-模当且仅当 A 是体（I, §9, No.~1, 第 115 页，定理 1）。这时单 A-模就是体 A 上的 1 维向量空间。

\begin{enumerate}
\def\labelenumi{\arabic{enumi})}
\setcounter{enumi}{1}
\item
  设 A 是一个主理想整环（VII, §1, No.~1, 第 1 页，定义 1）且不是体。对 A 的每个不可约元 \(\pi\)，A-模 \(A_s / (\pi)\) 都是单模，且每个单 A-模都同构于这样一个模（VII, §4, No.~8, 第 25 页，注 4）。对 \(n \geqslant 2\)，A-模 \(A_s / (\pi^n)\) 不可分解（VII, §4, No.~8, 第 24 页，命题 8）但不是单模。
\item
  设 K 是体，V 是体 K 上的一个非零右向量空间，A 是环 \(\operatorname{End}_{\mathrm{K}}(\mathrm{V})\) 的一个包含 V 的有限秩自同态的子环（例如 \(\mathrm{A} = \operatorname{End}_{\mathrm{K}}(\mathrm{V})\)）。我们来证明 V 是单 A-模：设 W 是 V 的一个非零 A-子模，x 是 W 的一个非零元素；存在 V 上的一个线性型 \(\varphi\)，使 \(\varphi(x) \neq 0\)（II, §7, No.~5, 第 300 页，定理 6）。对 V 中每个 y，映射 \(z \mapsto y\varphi(z)\)——它是秩 \(\leqslant 1\) 的线性映射——属于 A；于是我们有 Ax = V，从而更有 W = V，这证明 V 是单 A-模。
\end{enumerate}

命题 1. —— 设 A 是环。

\begin{enumerate}
\def\labelenumi{\alph{enumi})}
\item
  设 \(\mathfrak{m}\) 是 A 的一个左理想。A-模 \(\mathrm{A}_s / \mathfrak{m}\) 是单模当且仅当 \(\mathfrak{m}\) 是极大左理想。
\item
  设 M 是单 A-模，x 是 M 的一个非零元素。我们有等式 \(M = Ax\)，x 的零化子 m 是 A 的一个极大左理想，且映射 \(a \mapsto ax\) 通过过渡到商定义出从 \(A_{s}/m\) 到 M 的一个同构。
\item
  设 M 是一个非零 A-模。若对 M 的每个非零元素 x 都有 M = Ax，则 M 是单模。
\end{enumerate}

\(A_{s}/m\) 的子模都具有形式 n/m，其中 n 是 A 的包含 m 的左理想（I, §4, No.~6, 第 41 页，定理 4）；因此，A-模 \(A_{s}/m\) 是单模当且仅当我们有 \(m \neq A\)，且 A 的每个包含 m 的左理想 n 都满足 n/m = 0 或 \(n/m = A_{s}/m\)，即 n = m 或 n = A。这就证明了 a)。

在 b) 的假设下，Ax 是 M 的一个非零子模，从而等于 M。于是，映射 \(a \mapsto ax\) 通过过渡到商定义出从 \(A_{s}/m\) 到 M 的一个同构；因此 A-模 \(A_{s}/m\) 是单模，且由 a) 左理想 m 是极大的。这就证明了 b)。

在 c) 的假设下，设 N 是 M 的一个非零子模。若 x 是 N 的一个非零元素，则我们有 \(Ax \subset N\) 且 Ax = M，从而 M = N。于是 M 是单模。

推论 1. —— 若环 A 不化为 0，则存在单 A-模。

事实上，由克鲁尔定理（I, §8, No.~6, 第 104 页，定理 1），A 存在极大左理想。

推论 2. —— 设 A 是一个局部环（VIII，第 26 页，定义 1），r 是其极大理想。A-模 \(A_{s}/r\) 是单模，且每个单 A-模都同构于 \(A_{s}/r\) 。

注. —— 1) 设 A 是交换环，m 是 A 的一个理想。则 m 是 A-模 \(A_{s}/m\) 的零化子（II, §1, No.~12, 第 219 页，定义 11）。因此，若 m 与 \(m'\) 是 A 的不同理想，则 A-模 \(A_{s}/m\) 与 \(A_{s}/m'\) 不同构。存在忠实单 A-模（II, §1, No.~12, 第 219 页）当且仅当 (0) 是 A 的极大理想，即 A 是体。

\begin{enumerate}
\def\labelenumi{\arabic{enumi})}
\setcounter{enumi}{1}
\tightlist
\item
  我们可以给出一个非交换环 A 及其两个不同的极大左理想 m 与 \(m'\) 的例子，使得 A-模 \(A_{s}/m\) 与 \(A_{s}/m'\) 同构（VIII，第 52 页，习题 3）。
\end{enumerate}

